\begin{itemize}\setlength\itemsep{3pt}
  \item \textbf{Procedural input verification did not improve clinical
        tool-call accuracy.} Across four pre-registered primary contrasts
        (Bonferroni $\alpha=0.0125$), no verifier condition (E, D, C)
        significantly beat the no-verifier baseline B.
  \item \textbf{Under a tool-forcing system prompt the verifier
        \emph{actively hurt} accuracy by 9.0 pp} (McNemar $p < 10^{-15}$).
        118 vignettes flipped from correct to wrong; only 22 flipped the
        other way.
  \item \textbf{The harm came from the reference, not the verifier.}
        A single field (sex) accounted for $\sim$32\% of all flags
        under upfront full-schema extraction, despite being among the
        easiest fields to extract correctly --- implicating the
        extraction \emph{architecture}, not the verifier's comparison logic.
  \item \textbf{Reactive per-tool-call extraction \emph{partially}
        recovered the diagnostic markers but not accuracy.} Sex fell to
        13\% of flags, correct-rate on intervention rows doubled, and
        total token cost dropped $\sim$57\%, but the accuracy gain vs
        upfront B$'$+E was not significant ($+1.9$ pp, $p = 0.103$).
        Schema size is a contributor, not the sole bottleneck.
\end{itemize}
